\section{A valid family and its compact semidefinite formulation}
\label{family:section}

The example in \cref{family:counterexample} belongs to a parameter family
with the same five-edge geometry. The validity domain is larger than the
domain with five interior contacts. Retaining that larger domain is what
makes exact enforcement by a small semidefinite formulation possible.

For $h\in\R$ and $d_1,d_2,d_3,k\ge0$, set
$L=h-d_1x-d_2y+d_3z$ and $D=d_1+d_2-h$, and define
\begin{equation}
\label{family:definition}
q_{h,d,k}=L^2+2d_3kz(1-x-y)+k(2D+k)xy.
\end{equation}
In particular, $h$ is unrestricted in sign. Let $\mathcal F$ denote the
set of these polynomials, and let $\cone(\mathcal F)$ denote its conic
hull. Scaling all five parameters by $\alpha\ge0$ scales the polynomial
by $\alpha^2$.

\begin{theorem}[Validity]
\label{family:valid}
Every member of $\mathcal F$ belongs to $\mathcal P_3^+$. Hence every
inequality $\mathcal L_{m,Y}(q_{h,d,k})\ge0$ is valid on both
$\mathcal Q_3$ and $\mathcal H_3$.
\end{theorem}

\begin{proof}
The following polynomial identity holds for all real parameters:
\begin{equation}
\label{family:identity}
\begin{split}
q_{h,d,k}={}&(L-kxy)^2+2d_3kz(1-x)(1-y)\\
 &+k(2d_1+k)xy(1-x)+2kd_2xy(1-y)
   +k^2x^2y(1-y).
\end{split}
\end{equation}
To check it, expand the square and collect the remaining terms as
$2kxyL+2d_3kz(1-x-y)+k(2D+k)xy-k^2x^2y^2$.
After substituting $L$ and $D$, this is precisely the sum of the four
remaining terms in \eqref{family:identity}.
Each term on the right is nonnegative on $C_3$ under the stated parameter
restrictions. No sign condition on $h$ is used. The square coefficients
are $d_i^2$, so at an atom with diagonal surplus $s_i\ge0$ the evaluated
inequality has value
\[
q_{h,d,k}(x)+\sum_{i=1}^3d_i^2s_i\ge0.
\]
Convexity gives validity on $\mathcal Q_3$ and $\mathcal H_3$.
\end{proof}

Zero parameters are included in this result. For example, $k=0$ gives an
affine square, which already has a disjoint-support certificate. The
following theorem concerns the strict parameter region where the family
has five interior edge zeros.

\begin{theorem}[Exposed rays missed by disjoint certificates]
\label{family:exposed}
Assume
\begin{equation}
\label{family:strict-region}
h,d_1,d_2,d_3,k>0,\qquad h<\min\{d_1,d_2\},\qquad D+k<d_3.
\end{equation}
Then $q_{h,d,k}$ has exactly the five zeros
\begin{equation}
\label{family:contacts}
\begin{split}
a&=(h/d_1,0,0),\qquad b=(0,h/d_2,0),\\
c&=(1,0,(d_1-h)/d_3),\qquad
d=(0,1,(d_2-h)/d_3),\\
e&=(1,1,(D+k)/d_3).
\end{split}
\end{equation}
They lie in the relative interiors of their respective edges.
The ray $\R_+q_{h,d,k}$ is exposed in $\mathcal P_3$, and
$q_{h,d,k}\notin\mathcal D_3$.
\end{theorem}

\begin{proof}
The strict inequalities place each displayed zero in its edge interior.
The third and fourth terms of \eqref{family:identity} force any zero to
satisfy $xy=0$ or $x=y=1$. If $y=0$, the second term forces $z=0$ or
$x=1$; the square then gives $a$ or $c$. The case $x=0$ gives $b$ or $d$.
When $x=y=1$, the square gives $e$. Conversely all five points are zeros.

For exposedness, suppose that $r\in\mathcal P_3$ vanishes at these five
points. Write its coefficients in the convention of \cref{set:cones}:
\[
r=r_0+a_1x+a_2y+a_3z+b_1x^2+b_2y^2+b_3z^2
  +c_{12}xy+c_{13}xz+c_{23}yz.
\]
A univariate quadratic nonnegative on an interval and zero in its
interior is a nonnegative multiple of the square centered at that zero;
this includes the identically zero polynomial. Therefore
\[
r(x,0,0)=b_1(x-h/d_1)^2,\qquad
r(0,y,0)=b_2(y-h/d_2)^2,
\]
with $b_1,b_2\ge0$. Comparing constants shows that for some $\lambda\ge0$,
\begin{equation}
\label{family:contact-bottom}
b_1=\lambda d_1^2,\quad b_2=\lambda d_2^2,\quad
r_0=\lambda h^2,\quad
a_1=-2\lambda hd_1,\quad a_2=-2\lambda hd_2.
\end{equation}
The vertical restriction through $c$ is
$b_3(z-(d_1-h)/d_3)^2$. Evaluating at $z=0$ and using
$d_1-h>0$ gives $b_3=\lambda d_3^2$.
Vanishing of the three vertical derivatives at $c,d,e$ gives
\begin{align*}
a_3+c_{13}&=-2\lambda d_3(d_1-h),\\
a_3+c_{23}&=-2\lambda d_3(d_2-h),\\
a_3+c_{13}+c_{23}&=-2\lambda d_3(D+k).
\end{align*}
Solving yields
\begin{equation}
\label{family:contact-mixed}
a_3=2\lambda d_3(h+k),\qquad
c_{13}=-2\lambda d_3(d_1+k),\qquad
c_{23}=-2\lambda d_3(d_2+k).
\end{equation}
Finally, the restriction through $e$ gives
$r(1,1,0)=\lambda(D+k)^2$, and hence
\[
c_{12}=\lambda\bigl(2d_1d_2+k(2D+k)\bigr).
\]
Every coefficient of $r$ is now $\lambda$ times that of $q_{h,d,k}$.
The conclusion also holds when $\lambda=0$, in which case all coefficients
vanish. The functional
\[
r\longmapsto r(a)+r(b)+r(c)+r(d)+r(e)
\]
is nonnegative on $\mathcal P_3$ and has zero face exactly
$\R_+q_{h,d,k}$. It therefore exposes this ray.

For disjoint-certificate exclusion, use a term positive at the origin,
as in \cref{family:counterexample}. Write its affine factor as
$\alpha+\beta x+\gamma y+\delta z$, with $\alpha\ne0$ and
$A=\varnothing$. Its weight is positive at $a$ and $b$, so
$\beta=-\alpha d_1/h$ and $\gamma=-\alpha d_2/h$.
These nonzero coefficients force $B\subseteq\{z\}$. The weight is then
positive at $c$ and $e$. The zero at $c$ gives
$\delta=\alpha d_3/h$, but the affine factor has value $\alpha k/h\ne0$
at $e$, a contradiction.
\end{proof}

The exposed-ray conclusion is stronger than validity. A member in
\eqref{family:strict-region} cannot be decomposed into nonproportional
nonnegative quadratics. Its full mixed coefficients satisfy
$c_{12}>0$, $c_{13}<0$, and $c_{23}<0$. Their product is positive, which
agrees with the sign restriction in \cref{three:sign}.
For $(h,d_1,d_2,d_3,k)=(1/2,1,1,3,1)$, this theorem gives the polynomial
and zeros of \cref{family:counterexample}.

\subsection{Exact enforcement of all family inequalities}
\label{family:enforcement-subsection}

Let $\ell$ be a functional on quadratics, not necessarily a valid moment
functional. Write
$\ell(1)=t$, $\ell(x_i)=m_i$, and $\ell(x_ix_j)=Y_{ij}$, and form
\begin{equation}
\label{family:b-B}
b=\begin{pmatrix}-m_x\\-m_y\\m_z\\-Y_{xy}\end{pmatrix},\qquad
B=\begin{pmatrix}
Y_{xx}&Y_{xy}&-Y_{xz}&Y_{xy}\\
Y_{xy}&Y_{yy}&-Y_{yz}&Y_{xy}\\
-Y_{xz}&-Y_{yz}&Y_{zz}&m_z-Y_{xz}-Y_{yz}\\
Y_{xy}&Y_{xy}&m_z-Y_{xz}-Y_{yz}&Y_{xy}
\end{pmatrix}.
\end{equation}
Both depend linearly on the retained first and second moments.

\begin{theorem}[One order-five block and six auxiliaries]
\label{family:lmi}
For arbitrary $t,m,Y$, the inequalities
\[
\ell(q_{h,d,k})\ge0\qquad
(h\in\R,\ d_1,d_2,d_3,k\ge0)
\]
hold if and only if there is a symmetric matrix $N\in\Sbb^4$ satisfying
\begin{equation}
\label{family:homogeneous-lmi}
N_{ii}=0,\quad N_{ij}\ge0\ (i<j),\qquad
\begin{pmatrix}t&b^T\\ b&B-N\end{pmatrix}\succeq0.
\end{equation}
In particular, at $t=1$ they hold if and only if
$B-bb^T\in\COP_4$. The formulation uses one affine positive semidefinite
block of order five, six scalar nonnegative variables, and no moments of
degree greater than two.
\end{theorem}

\begin{proof}
Put $v=(d_1,d_2,d_3,k)^T$. Expanding
\eqref{family:definition} gives
\begin{equation}
\label{family:parameter-quadratic}
\ell(q_{h,d,k})=t h^2+2h b^Tv+v^TBv.
\end{equation}
At $v=0$, nonnegativity for every $h$ requires $t\ge0$.
If $t>0$, minimization over the unrestricted scalar $h$ gives
\[
\min_{h\in\R}\ell(q_{h,d,k})
  =v^T\bigl(B-t^{-1}bb^T\bigr)v,
\qquad h=-t^{-1}b^Tv.
\]
Thus the inequalities are equivalent to
$B-t^{-1}bb^T\in\COP_4$.

We use the classical order-four identity
\begin{equation}
\label{family:cop4}
\COP_4=\Sbb^4_++\mathcal N_4,
\qquad
\CP_4=\DNN_4,
\end{equation}
where $\mathcal N_4$ is the cone of symmetric entrywise nonnegative
matrices \citep{Diananda1962,MaxfieldMinc1962}.
It gives an exact decomposition
$B-t^{-1}bb^T=P+N$, with $P\succeq0$ and $N\ge0$ entrywise.
Moving the diagonal part of $N$ into $P$ permits $N$ to have zero diagonal.
The Schur complement of $t>0$ gives
\eqref{family:homogeneous-lmi}, and proves the converse as well.

If $t=0$, the expression in \eqref{family:parameter-quadratic} is
nonnegative for every real $h$ only if $b^Tv=0$ for all $v\ge0$,
or equivalently $b=0$. The remaining condition is $B\in\COP_4$.
By \eqref{family:cop4}, this again is exactly
\eqref{family:homogeneous-lmi}. Conversely, a positive semidefinite block
with zero top-left entry has $b=0$, so no additional case is needed.
This also covers the zero functional and all parameter degeneracies.
\end{proof}

The unrestricted sign of $h$ is essential to this proof. Requiring
$h\ge0$ would generally prevent elimination at $h=-b^Tv$ and would
produce a different condition. The nonnegative matrix $N$ is also
essential: setting $N=0$ can exclude valid cube moments. At the genuine
cube atom $(x,y,z)=(1/2,1/2,0)$, the matrix $B-bb^T$ has a zero first
diagonal entry but entry $(1,4)$ equal to $1/8$, so it is not positive
semidefinite. All family inequalities are nevertheless valid there.
The theorem is an SDP lift with auxiliary variables; it does not claim
an unlifted LMI or minimality of the stated lift size.

\subsection{Separation and extraction of a cut}
\label{family:separation-subsection}

At a normalized candidate moment point, let $A=B-bb^T$. The formulation
in \cref{family:lmi} can be tested as a feasibility SDP. Its equivalent
separation problem is
\begin{equation}
\label{family:separation-sdp}
\mu=\min\{\langle A,W\rangle: W\succeq0,\ W\ge0
             \text{ entrywise},\ \operatorname{tr}W=1\}.
\end{equation}

\begin{proposition}
\label{family:separation}
Problem \eqref{family:separation-sdp} has an attained optimum, and
\begin{equation}
\label{family:rankone-separation}
\mu=\min\{v^TAv:v\ge0,\ \|v\|_2=1\}.
\end{equation}
Thus $\mu<0$ if and only if some family inequality is violated.
A nonnegative vector $v$ with $v^TAv<0$, together with $h=-b^Tv$,
gives such an inequality.
\end{proposition}

\begin{proof}
The feasible set in \eqref{family:separation-sdp} is nonempty and compact.
For example, $(I+\mathbf1\mathbf1^T)/8$ is positive definite, strictly
positive entrywise, and has trace one. By $\DNN_4=\CP_4$, every feasible
$W$ has a finite decomposition $W=\sum_r u_ru_r^T$, with $u_r\ge0$.
Its trace condition gives $\sum_r\|u_r\|_2^2=1$. Discarding zero terms,
its objective is the convex combination
\[
\langle A,W\rangle
=\sum_r\|u_r\|_2^2
 \left(\frac{u_r}{\|u_r\|_2}\right)^T
 A\left(\frac{u_r}{\|u_r\|_2}\right).
\]
Conversely each normalized nonnegative rank-one matrix is feasible.
This proves \eqref{family:rankone-separation}; the last statement follows
from \eqref{family:parameter-quadratic} at $t=1$.
\end{proof}

A rank-one optimum exists, but an SDP solver need not return it.
A negative objective matrix is an aggregate certificate. Extracting
an individual cut from that matrix requires a completely positive
factorization, after which at least one factor has negative quadratic
value, or a separate search for a nonnegative vector. An exact
alternative using the $15$ nonempty coordinate supports is given in
\cref{comp:family-separation}. These statements concern exact arithmetic;
a numerical implementation must certify its proposed violation.

For the rational point in \cref{family:gap}, the parameters
$h=1/2$ and $v=(1,1,3,1)^T$ give $-1/40$. Keeping this $v$ fixed and
optimizing $h$ gives
\[
h=\frac{2361}{5000},\qquad
\ell(q_{h,d,k})=-\frac{644321}{25000000}.
\]
Hence the family block cuts off the point and strictly strengthens
each relaxation compared in \cref{family:comparison-subsection} when
added to that relaxation.

\subsection{Symmetry and selective use on larger problems}
\label{family:symmetries}

Coordinate permutations and complements preserve $C_3$, and therefore
preserve validity of the family. Their affine action on first and
second moments is given in \cref{set:symmetry}.
The diagonal surplus at an atom is unchanged, so the same transformations
preserve validity on $\mathcal H_3$ as well as on $\mathcal Q_3$.
Swapping $x$ and $y$ merely exchanges $d_1$ and $d_2$. Thus all symmetry
copies can be imposed using at most $24$ blocks: three choices of the
distinguished coordinate and eight complement patterns.

For a sparse quadratic model, one may add a family block on a selected
triple whose three edge moments are present. If an edge moment is absent,
introducing it is a further modeling choice. One may instead generate
individual inequalities by separation. The results above prove validity,
exact enforcement of the selected family, and strict improvement for
specified relaxation points. Completeness after every symmetry copy has
been included is the question examined in the next section.
